% Save as: module4-problem-set.tex
% ============================================================================
% Problem Set 4: Asymmetric Information, Moral Hazard & Adverse Selection
% Microeconomic Theory for Experimentalists & Applied Microeconomists
% Built from templates/problem_set_template.tex. Compile with pdflatex.
% ============================================================================
\documentclass[11pt]{article}
\usepackage{amsmath, amssymb, amsthm}
\usepackage[margin=1in]{geometry}
\usepackage{enumitem}

\newtheorem{question}{Question}

\title{Problem Set 4: Asymmetric Information, Moral Hazard \& Adverse
Selection\\
\large Microeconomic Theory for Experimentalists \& Applied Microeconomists}
\author{Due: \textbf{[date --- before Lecture 1 of Module 5]}}
\date{}

\begin{document}
\maketitle

\noindent\emph{Ground rules: collaboration on ideas is encouraged; write-ups
are individual. You may use AI tools for parts flagged with
$\langle$AI$\rangle$ and must attach the transcript. Show all calculations.}

\medskip

\noindent\textbf{Weights:} Part A --- 30 points. Part B --- 45 points.
Part C --- 25 points.

\bigskip

% ---------------------------------------------------------------------------
\section*{Part A: Theory --- Guided Calculations (30 points)}

\begin{question}[Piece rates, flat wages, and selling the job --- 15
points]
A worker produces output worth \$1 per unit; output equals effort $e$.
Effort costs the worker $c(e) = e^2/20$ dollars and is \emph{unobservable}
--- only output can be contracted on. The worker's outside option is 0.

\begin{enumerate}[label=(\alph*)]
  \item \textbf{(4 pts)} Under a piece rate $r$ per unit, the worker
        maximizes $r e - e^2/20$. Show that she chooses $e = 10r$, and
        compute her effort, earnings (gross pay), effort cost, and net
        payoff at $r = 0.4$.
  \item \textbf{(3 pts)} Under a flat wage (paid regardless of output),
        what effort does a purely self-interested worker choose, and why?
        One sentence on what this predicts about flat-wage workplaces ---
        a prediction Lecture 2's gift-exchange evidence complicates.
  \item \textbf{(5 pts)} The firm keeps output minus pay, so under piece
        rate $r$ its profit is $e - re = 10r(1 - r)$. Find the firm's
        profit-maximizing $r$, the resulting effort, firm profit, and
        worker net payoff. Then show that ``selling the job'' --- piece
        rate $r = 1$ plus a fixed franchise fee $F = 5$ --- achieves the
        first-best effort $e = 10$ and makes the firm strictly better
        off. \emph{Hint: first-best effort maximizes total surplus
        $e - e^2/20$.}
  \item \textbf{(3 pts)} Now output is noisy (equal to effort on average)
        and the worker is risk-averse. In words plus one inequality:
        why does the sell-the-job contract stop being optimal, and what
        does the optimal piece rate trade off?
\end{enumerate}
\end{question}

\begin{question}[Participation and incentive constraints --- 15 points]
A project yields revenue \$120 on success and \$30 on failure. With high
effort (personal cost \$12 to the agent) success probability is 0.9; with
low effort (cost \$0) it is 0.5. The agent's outside option is \$20.
Effort is invisible; the contract is a wage pair $(w_s, w_f)$ paid on the
outcome. Both parties are risk-neutral; wages may not be negative.

\begin{enumerate}[label=(\alph*)]
  \item \textbf{(4 pts)} Write the participation constraint and the
        incentive constraint --- each first in one plain-English
        sentence, then in one line of algebra. Simplify the incentive
        constraint to a statement about the spread $w_s - w_f$.
  \item \textbf{(5 pts)} Find the cheapest contract implementing high
        effort (both constraints bind). Report $(w_s, w_f)$, the expected
        wage, and verify that implementing high effort beats the firm's
        best flat-wage-low-effort alternative.
  \item \textbf{(3 pts)} Explain in two sentences why no spread larger
        than \$30 helps the firm --- and why, with any risk aversion at
        all on the agent's side, a larger spread strictly hurts. What
        economic role does each binding constraint play?
  \item \textbf{(3 pts)} If the agent were risk-averse, state (no
        computation) which constraint becomes more expensive to satisfy,
        what happens to the firm's cost of implementing high effort, and
        when the firm would give up and accept low effort.
\end{enumerate}
\end{question}

% ---------------------------------------------------------------------------
\section*{Part B: Experimental Design (45 points)}

\begin{question}[Decomposing the performance-pay effect]
Lazear (2000) found Safelite's switch to piece rates raised output
$\sim$44\%, roughly half from incentives (the same workers producing
more) and half from selection (more productive workers joining). Design
an experiment --- lab, online, or field --- that separately identifies
the incentive and selection effects of a piece-rate contract for a real
task. Specify:
\begin{enumerate}[label=(\alph*)]
  \item \textbf{(6 pts)} Subject pool, recruitment, and the real-effort
        task, with the properties the task must have (measurable output,
        a quality margin, no team production) and why each matters.
  \item \textbf{(12 pts)} Treatments. Your design must include (i) a
        stage or arm where workers are \emph{randomly assigned} to fixed
        wage vs.\ piece rate, and (ii) a stage or arm where workers
        \emph{choose} their contract before working. Say precisely which
        comparison identifies the incentive effect, which identifies
        selection, and what the contract-choice data themselves reveal.
  \item \textbf{(9 pts)} Measurement: output quantity AND quality (piece
        rates famously distort unmeasured margins), plus a baseline
        productivity measure taken \emph{before} any contract is
        revealed --- explain what the baseline buys you for the selection
        analysis.
  \item \textbf{(9 pts)} Hypotheses and rough sample-size reasoning:
        anchor effect sizes on Lazear's decomposition (total
        $\sim$44\%, roughly half incentives), state the per-arm $n$
        needed to distinguish the incentive effect from zero and from
        the total effect. Simulation-based reasoning is welcome;
        formulas are not required.
  \item \textbf{(9 pts)} The biggest confound in YOUR design and your
        answer to it.
\end{enumerate}
\end{question}

% ---------------------------------------------------------------------------
\section*{Part C: Silicon Pilot Interpretation $\langle$AI$\rangle$
(25 points)}

\begin{question}[The agents that matched both literatures]
You run \texttt{module4-simulation-lab.ipynb} with neutral personas. Your
silicon workers deliver effort $\approx 5$ under the \$2.50 piece rate
(the exact selfish-benchmark optimum), effort $\approx 2$ under the
stingy flat wage, and effort $\approx 7$ under the generous flat wage
(tight distributions --- within-cell standard deviations under 1) ---
a textbook incentive response AND a strong gift-exchange pattern,
simultaneously. Your lab partner celebrates: ``we've validated standard
theory and Fehr et al.\ in one run!''

\begin{enumerate}[label=(\alph*)]
  \item \textbf{(8 pts)} Explain why producing BOTH canonical patterns at
        once is itself grounds for suspicion rather than double
        validation, and give two mechanisms (at least one LLM-specific)
        that could generate this result.
  \item \textbf{(9 pts)} Design the diagnostic simulation that
        determines what is driving the flat-wage effort pattern. Your
        design must include (i) a \emph{parametric wage sweep} (many wage
        levels between \$10 and \$25, identical wording) and (ii) an
        \emph{intentions manipulation} (wage set by a computer vs.\
        chosen by the employer), and may add others (e.g.\ stripping the
        experiment framing). State the predicted result pattern under
        each mechanism from part (a). Run it if API budget allows and
        report; otherwise state the predictions precisely.
  \item \textbf{(8 pts)} A colleague wants to use this lab to pilot a
        real gift-exchange field experiment --- including choosing the
        generous wage level and predicting the effort response. State
        one thing the silicon pilot CAN safely inform, one thing it
        CANNOT (and why), and the design change this implies for the
        human study.
\end{enumerate}
Attach your AI transcripts and, if you ran code, the results CSV.
\end{question}

\end{document}
